% Stichworte:
% 17.4 Eigenwerte als Lösung der Bewegungsgleichung (Vorspann)
% - handeln statt andocken: "In diesem Abschnitt lösen wir ..."
% - Fahrplan: skalar -> System -> charakteristische Gleichung
%
% 17.4.1 Der Exponentialansatz im skalaren Fall
% - skalare Gleichung xdot = ax; Ableitung proportional zur Funktion
% - genau eine Funktion: e^{lambda t}; Lösung x0 e^{at}
% - R9 (x0 = Anfangszustand) und R3 (Tatsächlich ...) erfüllt
% - Rückverweis Buchteil (y' = -ky, Potenzreihenansatz)
%
% 17.4.2 Vom gekoppelten System zur Eigenwertgleichung
% - Hürde Kopplung an der echten pdot-Zeile
% - Wunschbild entkoppelt (Indikativ!); Weg: Eigentheorie
% - Ansatz einsetzen, kürzen -> Eigenwertgleichung; Taufe
% - "Man beachte": die Zeit ist verschwunden
%
% 17.4.3 Die Eigenwerte als Nullstellen der charakteristischen
%        Gleichung
% - Umformen -> homogenes System; Entweder-oder-Weiche (Evelyn S. 2)
% - triviale Lösung unbrauchbar -> det(A - lambda I) = 0
% - Polynom 4. Grades, vier Nullstellen, i. A. komplex
% - Schluss: "Das sind die vier Eigenwerte der Matrix A."

\section{Eigenwerte als Lösung der Bewegungsgleichung
\label{aircraft:section:4-eigenwerte}}
\kopfrechts{Eigenwerte}
In diesem Abschnitt lösen wir die Bewegungsgleichung.
Wir betrachten zuerst den einfachsten Fall, eine einzelne Gleichung
mit einer einzelnen gesuchten Funktion.
Danach übertragen wir den gefundenen Ansatz auf das System und
gelangen über die charakteristische Gleichung zu den Eigenwerten.

\subsection{Der Exponentialansatz im skalaren Fall
\label{aircraft:subsection:exponentialansatz}}
Wir betrachten die skalare Gleichung
\begin{equation}
    \dot{x}
    =
    ax
    \label{aircraft:eqn:skalar}
\end{equation}
mit einer gesuchten Funktion $x(t)$ und einer festen Zahl $a$.
Die Gleichung verlangt, dass die Ableitung proportional zur Funktion
ist.
Aus der Analysis kennen wir genau eine Funktion mit dieser
Eigenschaft, die Exponentialfunktion mit
\begin{equation}
    \frac{d}{dt} e^{\lambda t}
    =
    \lambda e^{\lambda t}.
    \label{aircraft:eqn:expableitung}
\end{equation}
Die Lösung der Gleichung~\eqref{aircraft:eqn:skalar} ist daher
\begin{equation}
    x(t)
    =
    x_0 e^{at}.
    \label{aircraft:eqn:skalarloesung}
\end{equation}
Die Konstante $x_0$ wird durch die Differentialgleichung nicht
festgelegt.
Sie ist der Anfangszustand, denn Einsetzen von $t = 0$ ergibt
$x(0) = x_0$.
Tatsächlich löst die Funktion~\eqref{aircraft:eqn:skalarloesung} die
Gleichung~\eqref{aircraft:eqn:skalar}, denn ihre Ableitung ist
$\dot{x}(t) = a x_0 e^{at} = a x(t)$.
In Kapitel~\ref{chapter:differentialgleichungen} wurde die verwandte Gleichung
$y' = -ky$ mit einem Potenzreihenansatz gelöst; auch dort war die
Lösung die Exponentialfunktion.

\subsection{Vom gekoppelten System zur Eigenwertgleichung
\label{aircraft:subsection:eigenwertgleichung}}
Das skalare Rezept überträgt sich nicht direkt auf die
Bewegungsgleichung, denn das System ist nach
Abschnitt~\ref{aircraft:section:2-dynamisches-system} gekoppelt.
Wünschenswert ist ein entkoppeltes System, in dem jede Gleichung nur
ihre eigene Funktion enthält.
Den Weg dorthin liefert die Eigentheorie der Matrix $A$.
Wir setzen dazu die Lösung mit der Exponentialfunktion an, in der
Form
\begin{equation}
    \vec{x}(t)
    =
    \vec{x}_0 e^{\lambda t}
    \label{aircraft:eqn:ansatz}
\end{equation}
mit einem festen Vektor $\vec{x}_0$.
Einsetzen in die Gleichung~\eqref{aircraft:eqn:bewegungsgesetz}
ergibt
\begin{equation}
    \lambda \vec{x}_0 e^{\lambda t}
    =
    A \vec{x}_0 e^{\lambda t}.
    \label{aircraft:eqn:eingesetzt}
\end{equation}
Da $e^{\lambda t} \neq 0$ ist, dürfen wir kürzen und erhalten die
\emph{Eigenwertgleichung}
\begin{equation}
    A \vec{x}_0
    =
    \lambda \vec{x}_0.
    \label{aircraft:eqn:eigenwertgleichung}
\end{equation}
Die Zahl $\lambda$ heisst \emph{Eigenwert}, der Vektor $\vec{x}_0$
\emph{Eigenvektor} der Matrix $A$.

\subsection{Die Eigenwerte als Nullstellen der charakteristischen
Gleichung
\label{aircraft:subsection:charakteristische-gleichung}}
Wir bringen beide Terme der Eigenwertgleichung auf eine Seite und
erhalten
\begin{equation}
(A - \lambda I)\vec{x}_0
=
\vec{0}
\label{aircraft:eqn:homogen}
\end{equation}
mit der Einheitsmatrix $I$.
Dies ist ein homogenes lineares Gleichungssystem für $\vec{x}_0$.
Es hat entweder genau eine Lösung, nämlich $\vec{x}_0 = \vec{0}$,
oder unendlich viele.
Die Lösung $\vec{x}_0 = \vec{0}$ ist unbrauchbar, denn sie beschreibt
keine Bewegung.
Unendlich viele Lösungen gibt es nur, wenn $A - \lambda I$ singulär
ist, also
\begin{equation}
    \det(A - \lambda I)
    =
    0
    \label{aircraft:eqn:charakteristisch}
\end{equation}
gilt, die \emph{charakteristische Gleichung}.
Ausgeschrieben ist sie ein Polynom vierten Grades in $\lambda$ und
hat vier Nullstellen, die im Allgemeinen komplex sind.
Das sind die vier Eigenwerte der Matrix $A$.